\question{What crosses the practice boundary, and what does it contain?}
\label{sec:payload}
\statusMeasured

\begin{shortanswer}
On the recommended configuration, \UplinkLogreg{} bits per practice per round: eleven coefficients
and an intercept. That vector is fitted to this practice's patients and is the object that
gradient-inversion research targets, so treating it as carrying information about those patients is
correct. A second, \emph{historical alternative} was materially worse: the tree-based protocol (FedXgbBagging,
evaluated on this report's pinned run and \textbf{deleted from the codebase on 2026-08-28} --- see
the scope note in the front matter) transmitted serialized trees whose split thresholds were
\emph{literal patient feature values} --- the highest-disclosure payload in the comparison and the
one secure aggregation cannot protect. What remains is the comparison between the two
logistic-regression protocols below.
\end{shortanswer}

\begin{table}[htbp]
  \centering
  \caption{What each protocol places on the wire, per practice per round. Measured in the test bed.}
  \begin{tabular}{lp{7.4cm}r}
    \toprule
    Protocol & What crosses the practice boundary & Bits/round \\
    \midrule
    FedProx & model coefficients --- a parameter vector fitted to this practice's patients & 352 \\
    FedMosaic & predictions + expertise on a \emph{public} cohort --- no parameter vector & 3,600 \\
    FedXgbBagging & serialized decision trees --- split thresholds are literal patient values & 21,521 \\
    \bottomrule
\end{tabular}

  \label{tab:payload}
\end{table}

\begin{figure}[htbp]
  \centering
  \includegraphics[width=0.82\textwidth]{payload}
  \caption{Uplink volume per practice per round, log scale, for the two logistic-regression
  protocols. (The pinned run also measured a tree-based protocol at roughly sixty times the
  coefficient vector; that series is historical --- the model was removed from the codebase on
  2026-08-28.)}
  \label{fig:payload}
\end{figure}

\subsection{Why the payload's \emph{shape} matters, not just its size}

The three protocols differ in kind, not only in degree:

\begin{itemize}
  \item \textbf{Coefficients} (FedProx, FedAvg) are an aggregate summary: each number is a
        partial derivative averaged over hundreds of patients. Inverting them to recover an
        individual is the subject of the gradient-leakage literature \cite{zhu2019leakage}.
  \item \textbf{Serialized decision trees} (FedXgbBagging -- \emph{historical}: this protocol was
        evaluated and removed from the codebase on 2026-08-28) are categorically different. A split
        threshold such as \emph{age $\le$ 67.5} is derived directly from patient values present in
        that practice's data. There is no averaging step to hide behind.
  \item \textbf{Predictions on a public cohort} (FedMosaic \cite{fedmosaic}) are different again.
        The practice transmits opinions about patients who are \emph{already public}, plus an
        expertise weight. No parameter vector exists for the server to invert.
\end{itemize}

\subsection{How much can a coefficient vector reveal?}

The gradient-inversion results counsel's assessment alludes to were established against deep
networks with millions of parameters, where a single update can carry enough information to
reconstruct an input image. The object here is different by orders of magnitude: \textbf{eleven
parameters}, averaged over a practice's entire training set, in a \emph{convex} optimisation whose
optimum is a property of the dataset as a whole rather than of any record in it.

That is an argument for lower risk, and it should be stated as an argument not as a conclusion.
It is exactly the kind of claim EDPB Opinion 28/2024 \cite{edpb2024ai} declines to accept without
testing, which is why \S\ref{sec:anontest} tests it rather than asserting it. The measured result
there is consistent with this reasoning, and its limits are stated with it.
